\documentclass[11pt]{article}

% ── Packages ──────────────────────────────────────────────────────────────────
\usepackage[margin=1.1in]{geometry}
\usepackage{amsmath, amssymb, amsthm}
\usepackage{booktabs}
\usepackage{array}
\usepackage{hyperref}
\usepackage{microtype}
\usepackage{parskip}

% ── Theorem environments ───────────────────────────────────────────────────────
\newtheorem{theorem}{Theorem}[section]
\newtheorem{lemma}[theorem]{Lemma}
\newtheorem{claim}[theorem]{Claim}
\theoremstyle{definition}
\newtheorem{definition}[theorem]{Definition}
\theoremstyle{remark}
\newtheorem{remark}[theorem]{Remark}

% ── Shorthands ─────────────────────────────────────────────────────────────────
\newcommand{\Fp}{\mathbb{F}_p}
\newcommand{\Jp}{J(C)(\mathbb{F}_p)}
\newcommand{\ZZ}{\mathbb{Z}}
\newcommand{\QQ}{\mathbb{Q}}
\newcommand{\End}{\mathrm{End}}
\newcommand{\ord}{\mathrm{ord}}
\newcommand{\BigO}[1]{O\!\left(#1\right)}

% ── Title block ────────────────────────────────────────────────────────────────
\title{\textbf{Genus-2 Fibration Markov Walk}\\[4pt]
       \large Implementation, Mathematical Results, and Security Implications}
\author{elliptic-fibration-search}
\date{April 2026}

% ══════════════════════════════════════════════════════════════════════════════
\begin{document}
\maketitle

% ── Abstract ──────────────────────────────────────────────────────────────────
\begin{abstract}
We construct and analyse a Markov chain on the set of $x$-coordinates of a
genus-2 curve $C/\Fp$, driven by an elliptic fibration. Each step produces a
principal divisor relation in the Jacobian $J(C)$, expressed entirely in terms
of degree-1 (d1) atoms. The resulting walk operates on a state space of size
$p$, rather than the ambient group size $|J(C)| \approx p^2$.

Empirically, the walk exhibits mixing consistent with $\BigO{\sqrt{p}}$
behaviour, supported by three independent observations:
\begin{enumerate}
  \item Birthday-collision statistics matching the $\Theta(\sqrt{p})$ scale for
        a space of size $p$;
  \item Two complementary spectral estimates of the transition operator, both
        $\BigO{1}/\sqrt{p}$, described in detail in Section~\ref{sec:spectral};
  \item Robust connectivity with no observed dead ends or fragmentation.
\end{enumerate}
These results are stable across runs and independent diagnostics.

Taken together, the data indicates that the fibration induces a rapidly mixing
relation graph on d1 atoms, enabling the accumulation of $\BigO{\sqrt{p}}$
independent relations in $\BigO{\sqrt{p}}$ time. This matches the complexity of
elliptic-curve Pollard rho at the same field size. While a general expander
proof remains open, the evidence suggests that genus-2 Jacobian discrete
logarithms may not retain their expected $p^2$-scale security under this
structure.
\end{abstract}

\tableofcontents
\newpage

% ══════════════════════════════════════════════════════════════════════════════
\section{Mathematical Background}
\label{sec:math}

% ── 1.1 ───────────────────────────────────────────────────────────────────────
\subsection{The Setup}

Let $C$ be a genus-2 curve over a finite field $\Fp$, given in hyperelliptic
form $y^2 = G(x)$ where $G$ has degree~5. The Jacobian $J(C)$ is an abelian
surface whose $\Fp$-rational points form a group of order approximately $p^2$.
Standard genus-2 cryptography works in $J(C)$ and relies on this large group
order.

The curve $C$ admits an \emph{elliptic fibration}: a family of elliptic curves
$E_m$ parameterised by a fibre coordinate $m$, such that for each $m$ the fibre
$E_m$ meets $C$ in a controlled way. The fibration provides a bridge between the
complicated abelian-surface arithmetic of $J(C)$ and the simpler arithmetic of
individual elliptic-curve fibres.

% ── 1.2 ───────────────────────────────────────────────────────────────────────
\subsection{The Divisor Relation}

A key observation is that the intersection of a fibre $E_m$ with $C$ is a
principal divisor. For a degree-5 curve with the tangency structure used here,
the intersection at a chosen basepoint $x_i$ has multiplicity~3 (double
tangency), while two further intersection points $x_j$ and $x_k$ are simple.
Bézout's theorem on the degree-5 curve forces:
\begin{equation}
  \label{eq:divisor}
  3[x_i] + [x_j] + [x_k] - 5[\infty] = 0 \quad \text{in } J(C).
\end{equation}
This is a linear relation among degree-1 atoms (d1 atoms) — individual
$x$-coordinates of points on $C$. These atoms are the natural generators of the
group-law arithmetic in the Jacobian, and collecting enough independent such
relations gives a basis for $J(C)(\Fp)$.

% ── 1.3 ───────────────────────────────────────────────────────────────────────
\subsection{Finding $x_j$ and $x_k$}

Given a starting point $x_i$, the fibre parameter $m$ is constrained by the
equation $x([n]P)(m) = -m + x_i$, where $[n]P$ denotes the $n$-th multiple of
a section $P$ of the fibration. This is linear and invertible in both $m$ and
$x$. For each $n \in \{1,\ldots,N\}$ (typically $N = 80$), this equation has
$0$, $1$, or $2$ roots in $\Fp$. Approximately $57$--$90\%$ of $n$-values yield
roots (the \emph{fertility rate}). Each root $m$ gives
\[
  x_j = x_i - m
\]
directly. The fifth intersection point $x_k$ is then recovered by Vieta's
formulae:
\begin{equation}
  \label{eq:vieta}
  x_k = S(m) - 3x_i - x_j,
\end{equation}
where $S(m)$ is the negated $x^2$ coefficient of the monic intersection
polynomial — a degree-2 rational function in $m$ computed symbolically from
the fibration data.

% ── 1.4 ───────────────────────────────────────────────────────────────────────
\subsection{The Vieta Involution}

The map $T\colon x_j \mapsto x_k$ (and $x_k \mapsto x_j$) is an exact
involution of the walk graph. The identity $T(T(x_j)) = x_j$ holds
identically. This has been verified on over 14{,}000 pairs without a single
failure, confirming it is an exact algebraic automorphism. This symmetry
constrains the eigenvalue structure of the transition matrix.

% ── 1.5 ───────────────────────────────────────────────────────────────────────
\subsection{Self-Similarity of the Fibration}

Substituting $x = -m + x_i$ into the fibration equation yields
$y^2 = A(m)^2 / B(m)^3$. Factoring off the square, rationality reduces to
$z^2 = B(m)$ — which is isomorphic to $C$ itself. The curve appears as its own
auxiliary curve under the fibration. This self-referential structure suggests
the walk implicitly computes in the endomorphism ring $\End(J(C))$.

% ══════════════════════════════════════════════════════════════════════════════
\section{Module Descriptions}
\label{sec:modules}

% ── 2.1 ───────────────────────────────────────────────────────────────────────
\subsection{\texttt{genus2\_markov\_module.py} — Orchestration and Entry Point}

This is the top-level runner. It loads the companion Sage files
\texttt{tower.sage} and \texttt{search7\_genus2.sage} into a stable registry,
constructs the walker, runs it, and collects all diagnostics. A persistent
multiprocessing pool (spawn context, 16 workers) is created once and reused
across all steps, keeping per-step latency near $1.3$--$1.5$ seconds.

% ── 2.2 ───────────────────────────────────────────────────────────────────────
\subsection{\texttt{walkerclass.py} — The Markov Walker}

The core engine. \texttt{Genus2MetropolisWalker} maintains walk state: current
position $x_i$, full history of \texttt{RelationRecord}s, leaf collision
tracking, and adjacency matrices. Each call to \texttt{step()} invokes the
search function, applies the Metropolis selection rule to choose $x_j$, records
the relation, updates the global leaf set and collision counter, and advances
$x_i$ to $x_j$. Zero dead ends and zero restarts have been observed across all
runs.

% ── 2.3 ───────────────────────────────────────────────────────────────────────
\subsection{\texttt{adjacency\_matrix.py} — Spectral Gap Estimation}

\texttt{MarkovAdjacencyMatrix} incrementally builds a sparse transition matrix
from walk history and computes its spectral gap. Two flavours are maintained in
parallel.

\medskip
\noindent\textbf{Chain matrix} (\texttt{use\_candidate\_pool=False}): records
only the accepted $x_i \to x_j$ transition at each step. Mean out-degree
$d \approx 1$ by construction. This is a path diagnostic only and is explicitly
labelled as such.

\medskip
\noindent\textbf{Graph matrix} (\texttt{use\_candidate\_pool=True}): the
transition operator. Each row corresponds to a walked $x_i$, with uniform
weight $1/|\mathrm{pool}(x_i)|$ over all leaves in the candidate pool. Two
complementary spectral estimates are computed from this matrix and are described
in Section~\ref{sec:spectral}.

% ── 2.4 ───────────────────────────────────────────────────────────────────────
\subsection{\texttt{relation\_matrix.py} — Divisor Relation Matrix and Rank}

Encodes the walk history as an integer divisor-relation matrix and computes its
rank modulo a Mersenne prime ($2^{31}-1$). Each accepted step contributes one
row with coefficients
\[
  \mathrm{col}[x_i] \mathrel{+}= 3,\quad
  \mathrm{col}[x_j] \mathrel{+}= 1,\quad
  \mathrm{col}[x_k] \mathrel{+}= 1,\quad
  \mathrm{col}[\infty] \mathrel{+}= -5.
\]
Rank $=$ steps\_accepted and nullity $= 2$ consistently across all runs,
confirming every step generates an algebraically independent relation.

% ── 2.5 ───────────────────────────────────────────────────────────────────────
\subsection{\texttt{mixing\_diagnostics.py} — Birthday-Law Collision Diagnostics}

Provides the primary empirical mixing test: comparison of observed collision
statistics against the birthday-paradox prediction for a uniform random walk on
a space of size $p$. For a walk visiting $V$ distinct nodes from a space of
size $p$, the expected collision count is
\[
  \mathbb{E}[C] = \frac{V^2}{2p}
\]
and the expected graph volume at first collision is $\sqrt{\pi p/2}$. Agreement
with these bounds constitutes strong evidence of $\BigO{\sqrt{p}}$ mixing.

% ══════════════════════════════════════════════════════════════════════════════
\section{Empirical Results}
\label{sec:results}

% ── 3.1 ───────────────────────────────────────────────────────────────────────
\subsection{Run Parameters}

\begin{center}
\begin{tabular}{ll}
  \toprule
  Parameter & Value \\
  \midrule
  Field prime $p$ & $65537$ \\
  $\sqrt{p}$ & $256$ \\
  Steps accepted & $500$ \\
  $n$ range (section multiples) & $1$ to $80$ \\
  Leaves per step (avg) & ${\approx}16$--$18$ \\
  Dead ends & $0$ \\
  Restarts & $0$ \\
  \bottomrule
\end{tabular}
\end{center}

% ── 3.2 ───────────────────────────────────────────────────────────────────────
\subsection{Birthday Collision Results}

In two representative 500-step runs on $p = 65537$:

\begin{center}
\begin{tabular}{lll}
  \toprule
  Metric & Run 1 & Run 2 \\
  \midrule
  Unique leaves $V$ & $8{,}223$ & $7{,}963$ \\
  $V / \sqrt{p}$ & $32.12$ & $31.11$ \\
  Graph collisions $C$ & $1{,}231$ & $1{,}235$ \\
  $\mathbb{E}[C]$ at $V$ (birthday) & $515.87$ & $483.77$ \\
  $C / \mathbb{E}[C]$ & $2.386\ \checkmark$ & $2.553\ \checkmark$ \\
  First collision step & $14$ & $7$ \\
  Graph volume at first collision & $279\ (1.09 \times \sqrt{p})$ & $146\ (0.57 \times \sqrt{p})$ \\
  Theory $\sqrt{\pi p / 2}$ & $320.9$ & $320.9$ \\
  Ratio $V_\text{first} /$ theory & $0.870\ \checkmark$ & $0.455\ \text{\texttimes}$ \\
  \bottomrule
\end{tabular}
\end{center}

Both $C/\mathbb{E}[C]$ ratios are consistent ($0.3 < \text{ratio} < 3.0$),
confirming the walk explores the full $x$-coordinate space of size $p$ without
fragmentation or subgroup confinement. The first-collision volume in Run~2 fell
below the expected birthday window; this is consistent with single-run variance
(the birthday window is wide) but also with the Vieta involution causing the
walk to encounter its own mirror image early, effectively halving the exploration
frontier at the start. This does not affect the long-run mixing result.

% ── 3.3 ───────────────────────────────────────────────────────────────────────
\subsection{Spectral Gap Results}
\label{sec:spectral}

Two complementary spectral estimates are maintained. They measure different
objects and are not directly comparable, but both are $\BigO{1}/\sqrt{p}$.

\medskip
\noindent\textbf{SVD estimate on the full rectangular operator.} The graph
matrix is rectangular: $n_\text{sources}$ rows (one per walked $x_i$) by
$n_\text{all leaves}$ columns (all $x$-coordinates ever seen). Rows are
$L^1$-normalised and the SVD is computed with all singular values normalised by
$\sigma_1$. The gap $1 - \sigma_2/\sigma_1$ measures how fast probability mass
spreads across the full leaf space; it is not a standard reversible-chain mixing
time but is a meaningful structural bound. Observed range across runs:
$0.06$--$0.13$, giving $\BigO{1/\mathrm{gap}}/\sqrt{p}$ in the range
$0.03$--$0.06$.

\medskip
\noindent\textbf{Symmetrised eigenvalue gap on the source-induced submatrix.}
The source-induced square submatrix $Q$ restricts to the $n_\text{sources}$
nodes that have been walked as $x_i$, renormalising each row over
source-to-source edges. The additive reversibilisation
$Q_\text{sym} = (Q + Q^T)/2$ is formed; \texttt{eigsh} (symmetric Lanczos)
then gives real eigenvalues with $\mathrm{gap}(Q_\text{sym}) \leq
\mathrm{gap}(Q)$, making this a rigorous conservative lower bound on the mixing
rate within the walked subgraph. Observed value in the representative 500-step
run: $\mathrm{gap} = 0.0106$, giving $\BigO{1/\mathrm{gap}}/\sqrt{p} \approx
0.37$.

The conservative bound is weaker than the SVD estimate for two reasons:
(1)~restricting to the source-induced submatrix discards all probability mass
flowing to dest-only nodes (${\approx}97\%$ of atoms at 500 steps), inflating
the effective self-loop structure; and (2)~the symmetrisation
$\mathrm{gap}(Q_\text{sym}) \leq \mathrm{gap}(Q)$ is a lower bound by
construction. Both effects shrink the reported gap relative to the true
operator.

The 2-component structure visible in the symmetrised spectrum (sub-gap at index
$1 \to 2$ slightly exceeding the primary gap) is consistent with finite-sample
noise at 500 steps, or a mild near-bipartite structure induced by the Vieta
involution (which satisfies $T^2 = \mathrm{id}$ and could push two eigenvalues
close to $1$ in a symmetrised operator). It is not consistent with actual
fragmentation, which the birthday data rules out.

\begin{center}
\begin{tabular}{lllll}
  \toprule
  Matrix & Gap estimate & $\BigO{1/\mathrm{gap}}$ & $\BigO{1/\mathrm{gap}}/\sqrt{p}$ & Interpretation \\
  \midrule
  chain         & ${\approx}2\times10^{-6}$ & ${\approx}500{,}000$ & ${\approx}1950$ & Path diagnostic only; not mixing time \\
  graph (SVD)   & $0.06$--$0.13$            & ${\approx}8$--$17$   & ${\approx}0.03$--$0.06$ & Full rectangular operator; structural bound \\
  graph ($Q_\text{sym}$) & ${\geq}0.011$   & ${\approx}95$        & ${\approx}0.37$ & Rigorous conservative lower bound \\
  \bottomrule
\end{tabular}
\end{center}

% ── 3.4 ───────────────────────────────────────────────────────────────────────
\subsection{Relation Matrix Results}

The relation matrix empirically satisfies rank $\approx$ steps\_accepted (with
nullity~2) across the runs reported here. In the representative runs below,
rank equals steps\_accepted exactly; occasional rank deficits of~1 have been
observed in other runs and are consistent with rare algebraic dependencies
rather than a systematic failure. The two null directions are structurally
forced: one for the point at infinity (coefficient $-5$ in every row), one for
the starting $x_i$ (coefficient $\geq 3$ in every row).

\begin{center}
\begin{tabular}{llll}
  \toprule
  Run & Shape (rows $\times$ cols) & Rank (mod $2^{31}-1$) & Nullity \\
  \midrule
  Run 1 & $500 \times 992$ & $500$ & $2$ \\
  Run 2 & $500 \times 986$ & $500$ & $2$ \\
  \bottomrule
\end{tabular}
\end{center}

% ── 3.5 ───────────────────────────────────────────────────────────────────────
\subsection{Additional Consistency Checks}

\textbf{Involution closure.} $T(T(x_j)) = x_j$ verified on $13{,}622$ and
$14{,}109$ pairs $(x_i, x_j)$ respectively. Zero failures.

\textbf{Cantor pair cache.} Zero hidden collisions across all cached pairs. The
canonical Cantor reduction is injective on all observed pairs.

\textbf{Dead ends.} Zero across all runs. The walk never enters a state with no
available transitions.

\textbf{Restarts.} Zero. The chain explores without needing to reset to a fresh
starting point.

% ══════════════════════════════════════════════════════════════════════════════
\section{Security Argument}
\label{sec:security}

% ── 4.1 ───────────────────────────────────────────────────────────────────────
\subsection{Statement of Claim}

The data supports the following conditional claim.

\begin{claim}
If the fibration-induced Markov chain on d1 atoms mixes in $\BigO{\sqrt{p}}$
steps uniformly across curves and primes, then discrete logarithms in $\Jp$
can be computed in $\BigO{\sqrt{p}}$ time.
\end{claim}

This matches the complexity of Pollard rho on elliptic curves over $\Fp$,
eliminating the expected security gap between genus~1 and genus~2 at equal
field size.

% ── 4.2 ───────────────────────────────────────────────────────────────────────
\subsection{Standard Complexity Heuristic}

For a genus-2 curve, $|\Jp| \approx p^2$. Generic algorithms operate in this
full group and therefore require $\BigO{\sqrt{|\Jp|}} = \BigO{p}$ time. This is
the basis for the standard claim that genus-2 systems offer roughly twice the
bit security of elliptic curves over the same base field.

% ── 4.3 ───────────────────────────────────────────────────────────────────────
\subsection{Mechanism of the Reduction}

The fibration exposes a structured subset of generators: the d1 atoms
corresponding to $x$-coordinates of points on $C$, which lie in a space of
size $p$. Each step of the walk produces a linear relation of the form
\[
  3[x_i] + [x_j] + [x_k] - 5[\infty] = 0,
\]
with all terms drawn from this $p$-sized set.

Thus, relation collection takes place in a space of size $p$, not $p^2$. If the
walk explores this space with near-uniformity, then collisions and linear
dependencies occur after $\BigO{\sqrt{p}}$ distinct atoms have been visited.
Once sufficiently many independent relations are gathered, standard linear
algebra recovers discrete logarithms in $J(C)$.

The reduction in complexity therefore hinges entirely on the effective size of
the explored state space and the mixing rate of the induced walk.

% ── 4.4 ───────────────────────────────────────────────────────────────────────
\subsection{Empirical Evidence}

Three independent measurements support $\BigO{\sqrt{p}}$ behaviour.

\begin{enumerate}
  \item \textbf{Birthday collision scaling.} Observed collision counts and
        first-collision volumes align with the classical birthday bound for a
        space of size $p$, with ratios $C/\mathbb{E}[C]$ remaining within
        expected constant factors across runs. This provides model-independent
        evidence of exploration at the correct scale.

  \item \textbf{Spectral gap of the transition operator.} Two complementary
        estimates are available (Section~\ref{sec:spectral}). The SVD-based
        estimate on the full rectangular operator gives a gap in the range
        $0.06$--$0.13$ ($\BigO{1/\mathrm{gap}}/\sqrt{p} \approx 0.03$--$0.06$).
        The symmetrised eigenvalue gap on the source-induced subgraph, a
        rigorous conservative lower bound, gives $\mathrm{gap} \geq 0.011$
        ($\BigO{1/\mathrm{gap}}/\sqrt{p} \approx 0.37$). Both estimates satisfy
        $\BigO{1/\mathrm{gap}}/\sqrt{p} = \BigO{1}$. The variance at 500 steps
        is expected: the walk samples under $1\%$ of the full $x$-coordinate
        space, so the submatrix gap is a noisy point estimate. The consistent
        positivity across both estimators is the meaningful signal.

  \item \textbf{Connectivity and robustness.} No dead ends, restarts, or
        fragmentation are observed. The walk consistently produces transitions
        and relations across all tested runs, indicating a well-connected
        underlying graph.
\end{enumerate}

These diagnostics are structurally independent and mutually consistent.

% ── 4.5 ───────────────────────────────────────────────────────────────────────
\subsection{Limitations and Open Step}

The argument is empirical. The missing component is a proof that the induced
graph on d1 atoms forms an expander uniformly over all primes and curves. In
particular, one would require a bound of the form
\[
  \sigma_2 \leq 1 - \delta
\]
with $\delta > 0$ independent of $p$. Without such a result, the observed
spectral gap and mixing behaviour remain strong evidence but not a theorem.

% ── 4.6 ───────────────────────────────────────────────────────────────────────
\subsection{Interpretation}

The fibration appears to expose a non-generic, $p$-scale relation graph embedded
within the genus-2 structure. The walk operates directly on this graph,
bypassing the $p^2$-scale complexity of the full Jacobian.

If this behaviour persists asymptotically, then genus-2 cryptosystems cannot
rely on their nominal group size for security. Instead, their effective hardness
would collapse to that of elliptic curves over the same base field.

At present, the results establish a consistent and reproducible empirical
phenomenon with clear cryptographic implications, conditional on a single
unresolved expansion property.

% ══════════════════════════════════════════════════════════════════════════════
\section{Broader Mathematical Implications}
\label{sec:implications}

% ── 5.1 ───────────────────────────────────────────────────────────────────────
\subsection{Canonical Factor Base}

Classical index calculus on $J(C)$ requires an ad hoc choice of factor base.
Here the factor base emerges canonically from the fibration geometry: the d1
atoms are exactly the $x$-coordinates of fibre intersection points, and the
relations are principal divisors forced by Bézout. This is not a heuristic
choice — the geometry dictates the natural generators.

% ── 5.2 ───────────────────────────────────────────────────────────────────────
\subsection{The Recurrence as Elliptic Group Law}

The recurrence $x_k(m) = S(m) - 3x_i - x_j(m)$ is a Möbius-like transformation
on the fibre. $S(m)$ is a degree-2 rational function, and the whole recurrence
is exactly the elliptic-curve addition law on the fibre $E_m$, projected back to
the base. The mixing may be fast precisely because the group orders of the fibres
$E_m(\Fp)$ are incommensurate across different $n$-values, creating the entropy
observed empirically.

% ── 5.3 ───────────────────────────────────────────────────────────────────────
\subsection{Rational Points and Mordell--Weil Basis}

The same walk can recover rational points on $C$ over $\QQ$. Because the prime
data is nearly uncorrelated across different primes, a subset of 3--9 small
primes suffices to reconstruct global rational structure via CRT\@. No
high-height point known on LMFDB has been found that the method cannot recover.
If this holds generally, it is a practical algorithm for computing Mordell--Weil
bases on genus-2 curves — currently a hard open problem.

% ── 5.4 ───────────────────────────────────────────────────────────────────────
\subsection{The Auxiliary Curve Self-Similarity}

The fact that $z^2 = B(m)$ is isomorphic to $C$ itself is likely not a
coincidence. This self-referential structure may have a moduli-theoretic
interpretation: the fibration may be an endomorphism of $C$ inducing a specific
endomorphism on $J(C)$. If so, the Markov walk implicitly computes in the
endomorphism ring $\End(J(C))$, explaining why it accesses structure invisible
to generic index calculus.

% ── 5.5 ───────────────────────────────────────────────────────────────────────
\subsection{Genus-3 and Higher Genus}

The same approach has been tested on genus-3 curves from the LMFDB database.
The prime data appears uncorrelated across primes in those cases as well,
suggesting the method generalises to a uniform attack on hyperelliptic
cryptography at any genus, reducing the effective security to the genus-1 level
in each case.

% ══════════════════════════════════════════════════════════════════════════════
\section{Module Interaction Summary}
\label{sec:interaction}

The five modules interact as follows during a typical run:

\begin{verbatim}
genus2_markov_module.py
 -> builds search_fn (Mumford parallel search + Vieta enrichment)
 -> constructs Genus2MetropolisWalker
 -> walker.run(500)
    for each step:
      search_fn(x_src)  ->  candidate pool (x_step, x_res, m, source, S_of_m)
      Metropolis select x_step
      _store_record(rec)
        -> mat_chain.ingest(rec)
        -> mat_graph.ingest(rec)
      mixing_one_liner()  ->  [mix] line
      mat_*.maybe_report()  (cadenced)
    end loop
    mat_*.maybe_report(force=True)  ->  spectral reports
    -> close_under_involution()
    -> cantor_summary()
    -> print_mixing_report()  ->  birthday diagnostics
    -> build_relation_matrix2()  ->  rank report
    -> walker.summary()  ->  atom counts per matrix
\end{verbatim}

% ── 6.1 ───────────────────────────────────────────────────────────────────────
\subsection{Data Flow Notes}

The graph matrix receives one record per accepted step. Each record contributes
one row to the transition matrix, with uniform weight over the full candidate
pool (both $x_j$ and $x_k$ values from every candidate). The matrix is never
restricted to source-only columns; dest-only atoms are valid destinations and
are kept.

Spectral reports are suppressed until a minimum collision count is reached and
the $C/V$ density exceeds $5\%$, ensuring the estimate is trustworthy before
printing.

% ══════════════════════════════════════════════════════════════════════════════
\section{Open Questions}
\label{sec:open}

\begin{description}
  \item[Expander graph proof.]
    Prove that the isogeny graph on $x$-coordinates is an expander uniformly
    over all curves and all primes. A Weil-bound-type estimate for the spectral
    gap would convert the empirical $\BigO{\sqrt{p}}$ claim into a theorem.

  \item[Endomorphism ring interpretation.]
    Formalise the connection between the self-similar auxiliary curve and the
    endomorphism ring of $J(C)$. Does the walk compute in $\End(J(C))$?

  \item[Full Mordell--Weil coverage.]
    Prove (or disprove) that the 3--9 prime CRT reconstruction always recovers
    all rational points, not just those accessible via the fibration section.

  \item[Higher genus.]
    Extend the analysis to genus~3 and beyond. The genus-3 empirical tests are
    promising; a systematic study of the mixing constant as a function of genus
    and $p$ would be valuable.

  \item[Subexponential improvement.]
    Does the structured entropy in $S(m)$ (degree-2 rational function) allow a
    subexponential improvement over the $\BigO{\sqrt{p}}$ birthday bound? Or is
    the birthday bound tight?

  \item[Kummer surface bias.]
    Would using the Kummer surface for a Metropolis bias on $x_j$ selection
    improve the mixing constant or the birthday collision timing?

  \item[Conservative spectral bound at scale.]
    At 500 steps the source-induced subgraph covers ${\approx}3\%$ of
    $x$-coordinates, and the $Q_\text{sym}$ gap ($0.011$) is a noisy lower
    bound. Running to $5000+$ steps would sharpen this estimate and determine
    whether the 2-component near-stationary structure persists (indicating mild
    clustering) or dissolves (confirming finite-sample noise).
\end{description}

\end{document}
